\section{Introduction}

Solvable learning models increasingly exhibit sharp changes in recovery,
attention mechanism, memorization, policy support, and collective order.  The
central unresolved question is not whether such transitions can be drawn, but
whether a theory calibrated before realism is added predicts an unseen
boundary afterward.  We ask:
\begin{quote}
Can a small set of macroscopic variables calibrated in a solvable model
predict phase boundaries in a held-out learning system, and can prediction
failure identify the missing variable?
\end{quote}

This question changes phase continuation from a descriptive atlas into a
falsifiable transport problem.  For an edge $u\to v$ in an assumption graph,
we fit effective parameters only at $u$, predict $\widehat g_{{\rm c},v}$, and score the
prediction against an untouched confirmatory grid at $v$.  A failed bridge is
scientifically useful only when an augmented observable restores prediction
on the same holdout.

Our four primary claims are: predictive boundary transport, localization of
non-additive assumption interactions, recovery by a new macrovariable, and
compute-efficient intervention inside the predicted critical window.  The
large experiment registry, immutable artifacts, and plotting grammar are
supporting infrastructure rather than scientific claims.

\resultfigure{figures/figure1_protocol.pdf}{Predictive phase continuation.  A
solvable anchor calibrates effective variables; all boundary and intervention
scores are evaluated on a blinded node.}
